\section{Limitations and Opportunities for Future Research}

The model does not currently incorporate environmental carrying capacity, so projected populations are not checked against the forage, land, or water base that would ultimately constrain them. Existing carrying-capacity methods diverge substantially in both approach and result: zone-based partitioning of agro-ecological data, remote-sensing-derived stocking rates \citep{meshesha2019}, and field-sampled biomass estimates can each be applied to the same rangeland and yield different sustainable-offtake thresholds, with sensitivity analysis showing swings of up to 20 percentage points in the area classified as overutilised depending on the modelling assumptions used \citep{zandler2023}. Population models that do embed carrying capacity, such as the stochastic Boran cattle model of \citet{tuffa2017}, treat its estimation as a substantial undertaking in its own right rather than a parameter that can be added on; integrating it here, ideally partitioned by agro-ecological zone and value chain, is better scoped as an independent follow-on study.

Data availability also constrained several value chains. Camel, goat, and sheep milk are frequently reported in Kenyan statistics as a combined ``other milk'' category rather than disaggregated by species, so chain-specific dairy parameters could not always be pinned to a domestic figure. For parameters with little or no Kenyan data, values were instead drawn from studies conducted in regions with comparable climate and breeds, a substitution consistent with FAO's own finding that livestock production parameters vary widely across production systems and geographies in Sub-Saharan Africa \citep{otte2002}, and one that introduces uncertainty proportional to how closely the donor region matches Kenyan conditions.

A further limitation not covered by the data issues above is scope: because the model is first-principles and driven by biological transition rates, it does not represent exogenous shocks that can move herd trajectories independently of demography, chief among them disease outbreaks, market price swings, and policy changes. Structured population models elsewhere in the literature that couple demographic transitions to herder market decisions, such as \citet{tuffa2017}, show that offtake responses to these external drivers can materially alter population trajectories; their omission here means the projections should be read as a demographic baseline rather than a forecast robust to shocks.

Addressing these gaps points to five priorities for future work: (i) a dedicated carrying-capacity study using zone- and value-chain-specific partitioning; (ii) targeted primary data collection for the species and parameters most affected by aggregation or missing coverage; (iii) integration of economic analysis, including price and cost feedbacks; (iv) incorporation of exogenous drivers such as disease, markets, and policy, potentially via a hybrid first-principles/machine-learning extension; and (v) disaggregating each value chain into independent, and where relevant breed-specific, sub-models for finer-grained decision support.
